\PageHeading{02}{赛题对照与完成证据}{按手册要求组织说明，不以功能名称代替实际验收}
手册将自研 Agent 的业务理解、产品依据、平台变化和真实运行作为重点。下表直接使用复赛评分项及分值，给出对应实现和证据边界；\textbf{不据此推算自评分数}。
\begingroup\small
\begin{tabularx}{\linewidth}{P{31mm}P{10mm}Y Y}
\rowcolor{Forest}\color{white}\bfseries 评分项 & \color{white}\bfseries 分值 & \color{white}\bfseries 当前设计与证据 & \color{white}\bfseries 正式验收关注点\\
业务场景理解 & 20 & 8类建议场景、目标客户、购买阶段和内容目标；案例A/B对比。 & 同产品换场景，核对切入点、卖点顺序及配图方案。\\
\rowcolor{Mist}产品事实准确性 & 20 & Fact保存原文、条件、来源；型号/系统约束及独立审稿。 & 需用指定ima知识库重新证明；本地快照不能代表最新事实。\\
图文生成与图片引用 & 15 & 配图ID、说明、来源和位置随稿件保存，实际HTML/PNG输出。 & 现场抽查ima原图与型号；非ima图片按手册不能计该项得分。\\
\rowcolor{Mist}平台适配与软文质量 & 15 & 四渠道分别规定结构、语气、图片数量及CTA，提供三平台案例。 & 对比实际语义与表达，不能只看模板或字数。\\
自研Agent与提示词 & 15 & 策划、已确认大纲写作、独立审稿、有限返修；提示词随项目交付。 & 核查业务约束和工具调用；不是直接转发ima写作。\\
\rowcolor{Mist}功能完成度与体验 & 10 & 七步流程、版本确认、取消、状态反馈、预览及导出；89项回归。 & 统一任务下检查可用性、异常恢复与输出追溯。\\
合规与安全 & 5 & 禁止无依据价格/承诺；凭证保护、输入限制、来源检查和受控渲染。 & 统一风险样本与凭证展示检查；仍需跟进依赖审计。\\
\bottomrule
\end{tabularx}
\endgroup
\Sub{最低功能板块如何落到产品中}
\begin{tabularx}{\linewidth}{P{38mm}Y}
\toprule
账号与知识库 & 接入页显示配置/连接状态；验证账号下的指定共享库，不回显密钥。\\
任务输入与补充 & 产品、客户、场景、渠道必选；缺少必要任务字段返回补充问题。\\
检索、策略与图文生成 & 按型号取事实与图片，以平台策略策划大纲；确认后生成正文。\\
核验、修改与导出 & 事实/图片/合规问题可定位；修改大纲并重新确认，导出实际文件。\\
\bottomrule
\end{tabularx}
\begin{attention}
\Lead{材料要求}\quad 初赛：GitHub、魔搭和作品说明书；复赛增加真实运行视频，并要求不少于3组结果、至少2类场景与2个平台。本说明书提供案例结构与现有证据；魔搭页面、录屏及ima正式样例需在提交前补齐并确认评审期可访问。
\end{attention}
\SourceNote{赛题手册 v4，第2、5、8--9页。表内分值为赛事权重，非项目得分。}
